\chapter{Multimodal Android Malware Analysis with \hybroid{}}
\label{chap:hybroid}

% Drafted under T-020 from notes/chapter4-source-records/. Governing decisions:
% D-022 the released repository is out of evidence and the paper governs;
% D-023 the corrected CICAndMal2017 category counts; D-024; D-025 the individual
% contribution; D-026 the two published claims the chapter declines to repeat.
% Classifier and ROC figures restored under T-002 from Hybroid Section 3.5.

\section{Research Problem and Contribution}
\label{sec:hybroid:1}

This chapter answers RQ2, stated in Section~\ref{sec:intro:rqs}: \rqTwoText{} It takes up
gap~\gapref{2}, which records that hybrid systems are rarely measured against each modality on
its own, so that the gain belonging to the combination cannot be separated from the gain
belonging to either part \cite{onwuzurike2018family, lindorfer2015marvin}. The question is
comparative, and the shape of the chapter follows from that. An accuracy figure for the combined
system answers nothing on its own. What answers RQ2 is the difference between the combined
system and each of its halves, measured on the same samples under one protocol, and
Section~\ref{sec:hybroid:11} reports it.

\hybroid{} observes an application twice. It recovers the program structure from the
application package without running it, and it observes the traffic the application produces
while running on a physical handset. Both observations are reduced to vectors, the two vectors
form one input, and a supervised classifier decides whether the application is malicious and, if
it is, which category it belongs to. Figure~\ref{fig:hybroid:architecture} gives the arrangement.
The design is described in Sections~\ref{sec:hybroid:5} to~\ref{sec:hybroid:8} and the evaluation
in Sections~\ref{sec:hybroid:9} to~\ref{sec:hybroid:11}.

\begin{figure}
  \centering
  \includegraphics[width=\linewidth]{hybroid-architecture.pdf}
  \caption{The Hybroid architecture. Every application sample provides a pair of inputs: the
  application package and a capture of the traffic it generates~(1). Both are analysed in
  parallel~(2). Androguard decompiles the package into opcodes, basic blocks and a control-flow
  graph, which are embedded into a 64-dimensional code representation, while Argus turns the
  capture into bidirectional flows that are aggregated, normalised and reduced to
  \netfeatures{}~network-flow features. The two representations are concatenated into one joint
  vector per sample~(3), on which a gradient-boosting classifier is trained and
  cross-validated~(4) for binary detection and for categorisation into four malware
  categories~(5). Redrawn by the candidate from the architecture published in
  \cite{norouzian2021hybroid}.}
  \label{fig:hybroid:architecture}
\end{figure}

Three claims are made for the work, and each is stated here in the form the evidence supports
rather than in the form the source publication used. The first is the flow pipeline: network
traffic captured from a real device is reduced to a fixed feature vector per application, through
bidirectional flow records, aggregation, scaling and a selection stage that leaves
\netfeatures{} features. The second is the program-code pipeline: instruction embeddings learned
from a large opcode corpus are composed into method-level vectors, which carry an
application-level function-call graph that a graph embedding converts into a single vector. The
third is the comparison itself, which measures each modality alone and the two together, under
one split, for detection and for category classification. The third depends on the first two and
is what Section~\ref{sec:hybroid:11} rests on.

Two limits on those claims belong here rather than in the limitations section, because they
qualify what is being claimed rather than what was found. The source publication names three
components, \texttt{opcode2vec}, \texttt{function2vec} and \texttt{graph2vec}, of which the
second is never defined anywhere in it and the third names a published method that the work does
not use. This chapter therefore describes the three steps and does not name them. And the
comparison in Section~\ref{sec:hybroid:9} is against four reimplemented systems of unequal
weight, which that section states when it reports the numbers.

The work was published at ISC 2021 \cite{norouzian2021hybroid}, with the candidate first of four
authors and the paper recording that the first two share first authorship. The candidate's own
contribution is stated in Section~\ref{sec:intro:publications}, which the affidavit covers: the
pipeline as a whole and the network-traffic analysis within it are his, as is the framework
through which the dataset was analysed, and he carried out the experiments and wrote the paper.
His contribution to the treatment of the graph data is partial, and this chapter does not present
the program-code pipeline as his own design.

\section{Related Work}
\label{sec:hybroid:2}

Section~\ref{sec:background:android:analysis} distinguishes static, dynamic and hybrid analysis
and reports that roughly one work in ten among surveyed deep-learning defences used both
\cite{liu2023deeplearning}. Section~\ref{sec:background:graphs:pairwise} defines the graph
objects. Neither is repeated here. What follows places \hybroid{} among the systems it was
measured against and among the work that has appeared since.

Static detectors read the application package. Features are drawn from the manifest and from the
code, and the families in use by 2021 were permissions and intents
\cite{arora2014malware, zulkifli2018android}, API calls
\cite{aafer2013droidapiminer, peiravian2013machine}, and structures recovered from the bytecode
\cite{gascon2013structural, li2019android, chen2018android}. DREBIN is the reference point among
them and one of the baselines in Section~\ref{sec:hybroid:9}: it gathers as many such features as
it can, from permissions and hardware requests through API calls to network addresses found in
the code, and embeds them in one vector space in which a linear model can be inspected to say
which features drove a verdict \cite{arp2014drebin}. Its breadth is also where it is exposed,
since the manifest entries and the identifiers a static analyser reads can be rewritten without
altering what the application does \cite{comparetti2010identifying}, which
Section~\ref{sec:background:android:obfuscation} treats in general terms.

A second static family works below the level of identifiers, on the instruction stream itself.
Sequences of opcodes have been used directly as features, with the effectiveness of $n$-grams over
them measured against families of malicious applications \cite{canfora2015effectiveness}, and
convolutional networks have been trained over raw opcode sequences to learn the relevant patterns
rather than to enumerate them in advance \cite{mclaughlin2017deep}. That line is the immediate
precursor of the representation in Section~\ref{sec:hybroid:5}, and the difference is where the
sequence stops. A model over a flat opcode sequence sees the program as one string. The
representation used here keeps the opcode alphabet and then organises it, first into methods and
then into the graph of calls between them.

Graph-based detectors use the structure rather than the identifiers. Adagio hashes neighbourhood
subgraphs of the function-call graph into an explicit feature map, which lets a linear classifier
be trained over graph structure at low cost and lets a verdict be traced back to the code
structures that produced it; it reports detection of 89 per cent of its malware
\cite{gascon2013structural}. It is the second static baseline in Section~\ref{sec:hybroid:9}.
apk2vec learns application embeddings from several graph views at once
\cite{narayanan2018apk2vec}. MaMaDroid takes a different route through the same object,
abstracting each call target to its package or family before modelling transitions between them,
which buys stability against the appearance of unseen methods at the cost of the callee's
identity \cite{onwuzurike2019mamadroid}.

All of these build the object \hybroid{} builds, and what separates them is how a vertex is
described. Adagio's vertex labels are hand-built from categories of instruction, and so are the
attributed control-flow graphs used for similarity between binaries
\cite{xu2017neural, yan2019classifying}. \hybroid{} learns the labels instead
\cite{bengio2013representation}. That choice has since acquired support it did not have at the
time: a benchmark study of graph-based Android classifiers reports that they exceed 94 per cent
accuracy on standard corpora yet lose up to 45 per cent under distribution shift on unseen
variants of known families, and it names structure-only function-call graphs as an inherent
limitation of those benchmarks, because topology alone misses the semantic patterns that
generalisation needs \cite{tran2025quantifying}. That study is a preprint under review. For the
wider design space of graph representation learning applied to malware, one survey covers the
extraction and architecture choices in full \cite{bilot2024survey}.

Dynamic detectors observe behaviour. Two branches exist, one reading system-level activity such
as system calls and API usage, the other reading what crosses the network. \hybroid{} belongs to
the second, which is cheaper to collect than system-call tracing and which exposes the part of an
application's behaviour that a remote operator depends on, since exfiltration and command traffic
have to leave the device to be of use \cite{malik2016credroid, arora2014malware}. Work in this
branch has ranged from rule-based inspection of what an application transmits
\cite{malik2016credroid} to learned models over traffic representations
\cite{wang2017malware, marin2021deepmal}. CICAndMal2017 was built to support it, and its authors
classify Android malware from flow features alone \cite{lashkari2018toward}, which makes their
system the dynamic baseline in Section~\ref{sec:hybroid:9}; they later added API calls to the same
pipeline \cite{taheri2019extensible}.

Granularity is the choice that divides this branch, and it runs from raw packets, which deep
models have been trained on directly \cite{marin2021deepmal}, through individual flows to whole
sessions \cite{wang2017malware}. The choice matters
more than it appears to, because it fixes what a label can attach to.
Section~\ref{sec:hybroid:7} returns to it.

Hybrid systems combine the two, and a systematic review of that literature is available and
replaces an enumeration here \cite{jain2026hybrid}. Two
points from the surrounding literature bear directly on this chapter. Measured under one
protocol, static analysis has proved at least as effective as dynamic analysis and the
combination has measured best \cite{onwuzurike2018family}; that comparison has rarely been
repeated, which is gap~\gapref{2}. And the modality pair used here was taken up again four years
later by HGDetector, which builds a static function-call graph and a network-behaviour graph from
the traffic, fuses them through a graph embedding, and evaluates on the same corpus
\cite{feng2025hgdetector}. Its result is conditional, and the condition is the useful part: the
gain from fusion is reported as roughly 4 per cent where the traffic features already capture an
application's behaviour, and 21 to 26 per cent where they do not. Section~\ref{sec:hybroid:11}
returns to that.

\section{Observation and Threat Model}
\label{sec:hybroid:3}

What a detector can see bounds what it can claim, and Chapters~\ref{chap:hybroid}
and~\ref{chap:hgann} differ in exactly that respect. This section fixes the observation model so
that the two are not read as competing designs.

\hybroid{} observes two things. The first is the Dalvik bytecode recovered from the application
package by a static analyser, from which methods, their instructions and the calls between them
are taken. The second is the traffic the application produces while running on a physical
handset, captured as packets and reduced to bidirectional flow records. Both observations are
made once, offline, before any verdict is given.

Several things are not observed. System calls are outside the model, as is code loaded at
runtime, and so is native code reached through the Java Native Interface. The contents of
encrypted payloads are also outside it, and that limit is measured rather than assumed: across
the four malware categories, the share of flows using TLS runs from 0.00 to 10.89 per cent, which
Section~\ref{sec:hybroid:7} reports and Section~\ref{sec:hybroid:12} takes up as a threat to
validity. The position in which the system is deployed follows from this. It analyses a sample
together with a capture of its traffic, which places it in an analysis pipeline rather than on
the handset.

The adversary is non-adaptive. The malware predates the system by several years and was not
written against this detector, no attacker knowledge is modelled, and no evasion attempt is
made against the trained model. Grey box in the sense fixed in
Section~\ref{sec:background:adversarial} is therefore not claimed here, since that term denotes
a level of attacker knowledge and this evaluation models none.
Section~\ref{sec:hybroid:12} states what follows from that, and what the attacks published since
2025 imply for a pipeline built over a function-call graph.

\section{Dataset and Experimental Protocol}
\label{sec:hybroid:4}

The evaluation uses CICAndMal2017 \cite{lashkari2018toward}, which supplies application packages
together with the traffic each one produced. Benign applications were collected from Google Play
across 2015 to 2017 and malicious ones from VirusTotal and the Contagio archive. The traffic was
captured from three physical handsets rather than from an emulator, with scripted user
behaviour and live accounts on messaging and social platforms, and in three states: immediately
after installation, before a reboot, and after it. Real devices are the reason this corpus was
chosen, since emulator captures do not reproduce the network behaviour of a handset in use.

\begin{table}
  \caption{CICAndMal2017 source-corpus composition. The category counts are recomputed from the
  dataset family listing and differ from the published values of 124, 112, 109 and 101, which do
  not sum to the stated total of 426.}
  \label{tab:hybroid:dataset}
  \centering
  \begin{tabular}{lrrr}
    \toprule
    Class & Families & Samples & Share \\
    \midrule
    Benign        &     & 1,700 & 80.0\,\% \\
    \midrule
    Adware        & 10  &   104 &  4.9\,\% \\
    Ransomware    & 10  &   101 &  4.8\,\% \\
    Scareware     & 11  &   112 &  5.3\,\% \\
    SMS malware   & 11  &   109 &  5.1\,\% \\
    \midrule
    Malware       & 42  &   426 & 20.0\,\% \\
    Total         &     & 2,126 & 100.0\,\% \\
    \bottomrule
  \end{tabular}
\end{table}

Table~\ref{tab:hybroid:dataset} gives the source-corpus composition. The published category rows
sum to 446 against the stated 426, and their percentages sum to 101.0. The recount closes on 426
across 42 families. The 47 exclusions are not identified by class. The dataset documentation
reports 5,065 benign applications installed, and the 1,700 used here is the number of benign
captures the distribution provides.

The source corpus pairs 2,126 applications with their traffic. Both analyses excluded the 47
graph-extraction failures and used the remaining 2,079 applications.
Model selection and evaluation use five-fold cross-validation with the results averaged over the
folds. Implementation used Python with
scikit-learn, TensorFlow and Keras, on a server with 32 cores, 128\,GB of memory and a 16\,GB
accelerator. Accuracy, precision, recall, the $F_1$ score and the area under the ROC curve
are reported, with $F_1$ preferred under the class imbalance Table~\ref{tab:hybroid:dataset}
shows. The averaging convention behind those metrics is not stated in the source, which
Section~\ref{sec:hybroid:9} returns to.

A second and larger experiment exercised the program-code pipeline alone. It used 45,592
malicious and 90,313 benign applications drawn from AndroZoo, VirusShare and VirusTotal, which is
135,905 samples at a malicious share of 33.5 per cent, split at random into 80 per cent for
training and 20 per cent for testing, and it reached an accuracy of 95.0 per cent with an $F_1$
score of 96.0 per cent. It is reported here as a check that the program-code pipeline scales
beyond a corpus of a few thousand samples. It is not a time-aware evaluation, and
Section~\ref{sec:background:evaluation} sets out why that distinction matters
\cite{pendlebury2019tesseract}.

\section{Opcode and Method-Level Representation}
\label{sec:hybroid:5}

Dalvik instructions carry an operation and its operands. Only the operation is embedded here, and
the reason is that the operands are register references whose values depend on the runtime that
executes them \cite{aosp2026dalvikbytecode, aosp2026art}, so they cannot be enumerated across
applications in the way a vocabulary must be. The operation, by contrast, is drawn from a fixed
alphabet of 216 opcodes, and samples within one malware family that share a malicious pattern
share the operations that implement it.

The embedding is learned rather than assigned. Each opcode is mapped to a real-valued vector by
the skip-gram formulation of word2vec, in which the model is trained to predict the operations
surrounding a given one \cite{mikolov2013distributed}. Training used a corpus of 18,240,542
opcodes extracted from applications drawn from AndroZoo. Two opcodes that occur in similar
positions therefore receive similar vectors, so the representation encodes the company an
operation keeps rather than an arbitrary index assigned to it, which is what a one-hot encoding
over the same alphabet would give.

\begin{figure}
  \centering
  \includegraphics[width=0.72\linewidth]{hybroid-opcode-powerlaw.pdf}
  \caption{Opcode frequency across the pre-training corpus. The upper panel names the twenty most
  common operations, and the lower panel plots all 216 in rank order, where the shape of the decay
  is the property the embedding relies on. Reproduced from the candidate's own publication
  \cite{norouzian2021hybroid}, as is Figure~\ref{fig:hybroid:correlation}.}
  \label{fig:hybroid:powerlaw}
\end{figure}

Borrowing a technique from language modelling needs a reason beyond convenience, and the
distribution of the alphabet supplies it. Counting opcodes across the corpus shows the twenty
most frequent following a power law, as Figure~\ref{fig:hybroid:powerlaw} shows, which is the
property that holds for words in natural language, which is the setting the skip-gram objective
was built for \cite{mikolov2013distributed}. The alphabet also behaves better
here than in the setting the technique came from. Natural language has an open vocabulary and a
model trained on one corpus meets unseen words in another; the Dalvik instruction set is closed
at 216 operations for the format the corpus was drawn from, so an embedding trained once covers
the applications the analyser reads, and the unseen-token case that troubles language models does
not arise in the same way. This is what makes a transductive method acceptable in a
detector that must generalise, and Section~\ref{sec:hybroid:6} returns to the point.

Discarding the operands has a cost, and it is specific rather than general. Every call
instruction receives the same vector whatever it invokes, because the target sits in the operand
field \cite{haq2021survey}. This is worth stating plainly at the point where the representation is
built, because the object built next is a graph whose edges are exactly those call targets. The
vertex features and the edge set are therefore recovered by two passes over the bytecode that
share no information, and the vertex features are blind to the structure that the edges record.

Above the instruction, a method is represented by combining the vectors of the instructions it
contains. The source publication calls this a basic-block embedding and gives it as a weighted
mean, $\tilde{f} = \sum_{i} w_i x_i / \sum_i w_i$, where $x_i$ is the $l$-dimensional embedding of
the $i$-th opcode. Two things about that definition are carried forward as they stand rather than
tidied. The weights $w_i$ are not specified anywhere in the source, which leaves the mean
undefined. And the formula is written over the $n$ opcodes of a function while the prose around it
calls the result the embedding of a basic block, so the two do not describe the same object.
Section~\ref{sec:hybroid:12} records both.

What the vector is used for is stated unambiguously, and that is what
Section~\ref{sec:hybroid:6} needs. It becomes the feature attached to a vertex of the method's
control-flow graph, whose vertices are basic blocks.

\section{Program-Graph Representation}
\label{sec:hybroid:6}

Section~\ref{sec:background:graphs:pairwise} defines the control-flow graph, the function-call
graph, and the composition in which per-method embeddings become the vertex features of a
call graph. \hybroid{} builds the composition. Control-flow graphs are extracted per method with
Androguard \cite{desnos2012android}, each one is reduced to a single vector, and those vectors
describe the vertices of the application's function-call graph, which is the object the
classifier consumes. Both terms are needed, and they are not interchangeable.

\begin{figure}
  \centering
  \includegraphics[width=\linewidth]{hybroid-code2vec.pdf}
  \caption{The program-code branch. The static analyzer extracts the control-flow graph, opcodes
  and basic blocks from each APK. Opcodes are embedded first, every method is represented from
  the tagged control-flow graph of the basic blocks it contains, and a graph embedding reduces
  the resulting function-call graph to a single 64-dimensional feature vector. Redrawn by the
  candidate from the arrangement published in \cite{norouzian2021hybroid}.}
  \label{fig:hybroid:code2vec}
\end{figure}

Figure~\ref{fig:hybroid:code2vec} sets out the branch as a whole. Two properties of that object
bound what any detector built on it can claim, and both are set out
in Section~\ref{sec:background:graphs:pairwise}. A call graph recovered without running the
program over-approximates, because the analysis admits edges that no execution takes. It also
under-approximates, since calls made through reflection, through code loaded at runtime or through
native libraries leave no edge for a static analyser to find, and Section~\ref{sec:hybroid:3}
records that all three are outside what is observed here. The object is in that sense an artefact
of the analyser rather than a property of the program, since which edges are found depends on the
toolkit and on how the component lifecycle is modelled. Neither property is
peculiar to this work. Both mean that the structure being classified is a static estimate of the
program's shape, and Section~\ref{sec:hybroid:12} states what follows when an adversary is
allowed to influence that estimate.

Reducing the graph to a vector could be done by counting substructures, which is what the kernel
methods discussed in Section~\ref{sec:hybroid:2} do \cite{gascon2013structural}. The route taken
here instead propagates information along the edges and lets the features of a vertex be shaped
by its neighbourhood.

The reduction from a graph to a vector uses structure2vec \cite{dai2016discriminative}. A
$p$-dimensional vector $\mu_v$ is attached to each vertex $v$ and initialised at random, then
updated for $T$ rounds by a function of the vertex's own features and the sum of its neighbours'
current vectors,
%
\begin{equation}
  \mu_v^{(t+1)} = F\!\left(x_v, \sum_{u \in N(v)} \mu_u^{(t)}\right), \qquad \forall v \in V,
  \label{eq:hybroid:meanfield}
\end{equation}
%
with $F$ realised as
%
\begin{equation}
  F\!\left(x_v, \sum_{u \in N(v)} \mu_u^{(t)}\right)
  = \tanh\!\left(W_1 x_v + \sigma\!\left(\sum_{u \in N(v)} \mu_u^{(t)}\right)\right),
  \label{eq:hybroid:update}
\end{equation}
%
where $W_1$ is trainable and $\sigma$ is an $n$-layer fully connected network,
%
\begin{equation}
  \sigma(l) = P_1\,\mathrm{ReLU}\!\left(P_2 \cdots \mathrm{ReLU}(P_n l)\right).
  \label{eq:hybroid:sigma}
\end{equation}
%
After $T$ rounds the vertex vectors are averaged and transformed by a second trainable matrix,
which yields one vector per graph. This is the operator
Section~\ref{sec:background:learning:messagepassing} describes in words, instantiated on the
graphs built here. The number of rounds fixes how far information travels: after $T$ updates a
vertex has been reached by everything within $T$ steps of it, so $T$ sets the size of the
neighbourhood any single vertex vector summarises. Averaging at the end rather than summing keeps
the result comparable across graphs of different sizes, which matters when applications differ by
orders of magnitude in how many methods they contain.

One step in this pipeline is not specified in the source publication, and the chapter states so
rather than supplying a plausible reconstruction. Equations~\ref{eq:hybroid:meanfield}
to~\ref{eq:hybroid:sigma} are applied at the level of a method's control-flow graph, and the
averaging inside that procedure pools the basic blocks of one method. The step that combines the
resulting per-method vectors into the single 64-dimensional vector standing for the whole
application is named as the third step of the pipeline and is specified by no equation. Nothing in
the source refers to a caller or a callee in symbols, so the composition is stated there as a
design and is reproduced here as one.

A note on transduction closes the description. The opcode embedding is transductive, since
word2vec needs the whole alphabet at training time. That alphabet is the fixed Dalvik instruction
set, so the requirement is met by construction rather than by assumption about the corpus. The
graph embedding is inductive, and graphs in the test partition are unseen when it is trained.

\section{Network-Flow Representation}
\label{sec:hybroid:7}

\begin{figure}
  \centering
  \includegraphics[width=\linewidth]{hybroid-dynamic.pdf}
  \caption{The network branch. Argus reduces a raw capture to bidirectional flow events, the
  events belonging to one application are aggregated into a single record, every attribute is
  normalised to a common range, and feature selection retains the \netfeatures{}~most informative
  features. Redrawn by the candidate from the arrangement published in
  \cite{norouzian2021hybroid}.}
  \label{fig:hybroid:dynamic}
\end{figure}

Section~\ref{sec:background:networks:flows} sets out the flow abstraction.
Figure~\ref{fig:hybroid:dynamic} gives the branch. Here the captures are processed with Argus \cite{qosient2026argus}, which produces bidirectional records grouped by the
usual five-tuple of source and destination address, source and destination port, and protocol,
and reports roughly forty statistics for each record. Two of those statistics are categorical,
the direction of the transaction and the state flags, and both are one-hot encoded; the rest are
numeric.

An application produces many flows, and the representation reduces them to one record. The values
of each feature are averaged across all flow events belonging to one capture, which yields a
single feature vector per application and matches the granularity of the program-code vector so
that the two can be combined. Two consequences follow and neither is stated in the source
publication. The label attaches to the application rather than to any individual flow, which
avoids the assumption that every flow emitted by a malicious application is itself malicious.
Against that, averaging discards the order and the timing of the capture, and a small number of
malicious flows is averaged together with whatever background traffic the handset produced during
the same window.

The averaged features are scaled to the unit interval by min-max normalisation, on the ground
that features whose ranges differ by orders of magnitude would otherwise contribute unequally to
any distance computed over them. That rationale applies to the two learners which were tried and
dropped rather than to the three reported in Section~\ref{sec:hybroid:8}, since a tree is
unaffected by any monotone rescaling of a feature. The source publication states that standard
scaling was considered and rejected.

\begin{table}
  \caption{The \netfeatures{} network-flow features retained after selection, grouped by the kind
  of quantity each one measures.}
  \label{tab:hybroid:flowfeatures}
  \centering
  \begin{tabular}{ll}
    \toprule
    Kind & Features \\
    \midrule
    Volume and duration & mean record duration, total bytes, source bytes \\
    Transport state     & destination TCP window, source loss and retransmission flag, \\
                        & destination out-of-order flag, transaction direction \\
    Path and service    & source and destination TTL, source and destination TOS byte \\
    Connection identity & source and destination TCP base sequence number \\
    \bottomrule
  \end{tabular}
\end{table}

Selection reduced the roughly forty statistics to the \netfeatures{} listed in
Table~\ref{tab:hybroid:flowfeatures}, using Pearson correlation, an extremely randomised trees
estimator and a univariate test that keeps the highest-scoring features. Redundancy among the
survivors was checked with Kendall's correlation, shown in
Figure~\ref{fig:hybroid:correlation}. The grouping in the table is this chapter's
addition, and it is there because the four groups do not have equal standing as descriptions of
application behaviour. Volume, duration and transport state are produced by what the application
does. The remaining two groups are not: the TTL reflects the path and the operating system of the
endpoint, the TOS byte is set by few stacks and is commonly zero, and the TCP base sequence
numbers are required to be unpredictable per connection and so cannot carry behaviour at all.
Section~\ref{sec:hybroid:12} takes this up, together with a published audit of the same feature
family in a different corpus.

\begin{figure}
  \centering
  \includegraphics[width=0.82\linewidth]{hybroid-flow-correlation.pdf}
  \caption{Kendall correlation among the retained flow features, used to check that the selected
  features are not redundant with one another.}
  \label{fig:hybroid:correlation}
\end{figure}

A separate observation concerns what the traffic carried, and
Table~\ref{tab:hybroid:protocols} records it. Cleartext HTTP accounts for between 29.22 and 61.38
per cent of flows across the four malware categories, and TLS for between nothing at all and
10.89 per cent. The remainder, which runs from 27.73 to 70.78 per cent depending on the category, is neither named nor characterised in the source publication. \begin{table}
  \caption{Share of flows by protocol, per malware category. The two columns do not sum to 100
  per cent and the source publication does not say what the remainder is.}
  \label{tab:hybroid:protocols}
  \centering
  \begin{tabular}{lrrr}
    \toprule
    Category & HTTP & TLS & Unaccounted \\
    \midrule
    Adware      & 52.00\,\% &  8.00\,\% & 40.00\,\% \\
    Ransomware  & 29.22\,\% &  0.00\,\% & 70.78\,\% \\
    Scareware   & 61.38\,\% & 10.89\,\% & 27.73\,\% \\
    SMS malware & 52.20\,\% & 10.28\,\% & 37.52\,\% \\
    \bottomrule
  \end{tabular}
\end{table}

Flow-level features
were chosen over inspection of payloads for reasons of privacy and cost, and that choice has
become the only available one since: beyond transport encryption, a substantial share of
applications encrypt content themselves before sending it, with one measurement of 12,598
applications finding 22.92 per cent doing so and most of those transmitting device-identifying
information \cite{pourali2022hidden}.

\section{Multimodal Fusion and Classification}
\label{sec:hybroid:8}

Two classifiers appear in this pipeline and they do different work, which is worth separating
before either is described, because a reader who meets both without warning will take one of them
for a mistake.

The first sits at the end of the program-code branch. A two-layer fully connected network takes
the graph vector and is trained to discriminate, with malicious samples labelled $+1$, benign
samples $-1$, and a decision taken at zero. Its purpose is to fit the parameters of the
representation, since the matrices in Equations~\ref{eq:hybroid:meanfield}
to~\ref{eq:hybroid:sigma} have to be learned against some objective and detection is the
objective available. What leaves the branch is the 64-dimensional vector, and the network that
produced it plays no further part.

One consequence deserves stating, because it bears on Section~\ref{sec:hybroid:10}. The objective
that fits the program-code representation is binary, and the same representation is then used for
the five-class task. A representation tuned to separate malicious from benign need not be equally
informative about which kind of malicious an application is, and nothing in the source
publication indicates that the embedding was refitted for the second task. Whether the
categorisation results would improve under a representation fitted to that objective was not
measured.

The second classifier acts on both modalities. The 64-dimensional program-code vector and the
\netfeatures{} flow features are taken together as one input, on which a supervised learner is
trained for detection and, separately, for category classification. The source publication
describes the two groups as forming a single input vector and does not name the operator that
combines them, so this chapter does not name it either. The proportions of that input are however
fixed by the two branch designs and they are uneven: 64 of the 77 dimensions come from the
program code, so the static half supplies roughly five dimensions in every six. No weighting is
applied to correct for this, and none is reported as having been considered.

Five learners were tried. Support vector machines and naive Bayes performed worst and were
dropped, leaving a decision tree, a random forest and a gradient boosting classifier, which are
the three reported throughout Sections~\ref{sec:hybroid:9} to~\ref{sec:hybroid:11}. No figures
are given for the two that were dropped, so the margin by which they were rejected is not on
record. The three that remain are tree-based learners of increasing capacity, which is convenient for
the comparison in Section~\ref{sec:hybroid:11}: the decision tree fits axis-aligned splits on a
single model, the random forest averages over many such models, and gradient boosting fits them
in sequence against the residual of what came before. Reading the three together therefore says
something about how much capacity the fused representation needs, and not only about which
learner wins.

Combining the modalities at the level of features, rather than training a model on each and
combining their decisions, is a design choice, and it was the ordinary one at the time. Its
consequences were not visible in 2021. Section~\ref{sec:hybroid:11} reports what the measurements
here show and sets them beside the evidence that has accumulated since.

\section{Binary Detection Results}
\label{sec:hybroid:9}

\begin{table}
  \caption{Detection compared with four reimplemented systems, read from the published figures.
  The four systems differ considerably in weight: DREBIN is a full conference paper, Adagio and
  CIC 2017 are papers of ten and seven pages, and the SVM entry is a two-page poster. The reported
  performance values come from Figures~6 and~7 of \cite{norouzian2021hybroid}.}
  \label{tab:hybroid:baselines}
  \centering
  \begin{tabular}{lrrrrr}
    \toprule
    System & Accuracy & $F_1$ & Recall & Precision & AUC \\
    \midrule
    SVM \cite{dai2015svm}                & 0.939 & 0.860 & 0.856 & 0.865 & 0.876 \\
    DREBIN \cite{arp2014drebin}          & 0.954 & 0.867 & 0.873 & 0.861 & 0.961 \\
    Adagio \cite{gascon2013structural}   & 0.893 & 0.932 & 0.913 & 0.953 & 0.867 \\
    CIC 2017 \cite{lashkari2018toward}   & 0.876 & 0.872 & 0.878 & 0.871 & 0.914 \\
    \midrule
    \hybroid{}                           & 0.970 & 0.969 & 0.970 & 0.970 & 0.996 \\
    \bottomrule
  \end{tabular}
\end{table}

Table~\ref{tab:hybroid:baselines} reports the comparison published for binary detection
\cite{norouzian2021hybroid}. \hybroid{} has the highest point value in each column. Two
qualifications accompany the comparison because they affect how the values can be read.

The first concerns the comparison set. The four systems are not of equal weight, and the SVM
entry is a two-page poster \cite{dai2015svm}, so the phrase state of the art would carry more
than the table supports.

The second qualification concerns the metrics. For a binary problem
in which accuracy is computed over the whole test set and precision and recall over the malicious
class, the three quantities determine the malicious share $p$ of that set through
$\mathrm{acc} = 1 - p\left[(1-R) + R(1-P)/P\right]$. Solving for each row of
Table~\ref{tab:hybroid:baselines} gives 22.0 per cent for the SVM entry, 17.2 for DREBIN, and
81.0, 49.2 and 50.0 per cent for Adagio, CIC 2017 and \hybroid{}. Depending on the labels of the 47
exclusions, the evaluated malware share lies between 18.2 and 20.5 per cent. The last three values
remain incompatible with that interval. The source publication does not state which averaging
convention produced its figures, and no raw result survives against which to check. The chapter
therefore preserves the published values and does not treat their differences as estimates made
under one documented metric convention.

\begin{figure}[t]
  \centering
  \includegraphics[width=0.72\linewidth]{hybroid-detection-roc.pdf}
  \caption{Receiver operating characteristic curves for binary malware detection, reproduced
  from the published vector content of Figure~7 in \cite{norouzian2021hybroid}. The source labels the curves as macro
  averages and reports the AUC values in the legend. It provides neither a selected operating
  threshold nor fold-level variation.}
  \label{fig:hybroid:detection-roc}
\end{figure}

Figure~\ref{fig:hybroid:detection-roc} shows the corresponding threshold curves. \hybroid{} has
the largest reported AUC at 0.996, followed by DREBIN at 0.961. The dynamic CIC 2017 system has an
AUC of 0.914, ahead of the SVM and Adagio entries despite its lower accuracy in
Table~\ref{tab:hybroid:baselines}. The metric therefore changes the ordering of the baselines.
The AUCs summarise ranking across thresholds, while the paper reports no threshold at which the
detector would operate. It also supplies no fold-level curves or confidence intervals, so the
differences remain descriptive point estimates \cite{norouzian2021hybroid}.

Read by modality, three of the four baselines are static and one is dynamic. The comparison still
mixes the implementation choices of separate systems with the observation modality, so it cannot
isolate a modality effect. Section~\ref{sec:hybroid:11} uses the within-system comparison for that
purpose.

\begin{figure}[t]
  \centering
  \includegraphics[width=\linewidth]{hybroid-detection-classifiers.pdf}
  \caption{Binary detection by learner and observation modality, with each reported metric in its
  own panel. Every row carries one marker per modality, joined by a rule spanning the values it
  reports. The values are the five-fold means printed to two decimal places in Figure~8 of
  \cite{norouzian2021hybroid}, checked against the bar geometry of that figure. The score axis is
  truncated to 0.86--0.98, the interval those values occupy. No fold-level scores are published,
  so the figure carries no error bars. Redrawn by the candidate.}
  \label{fig:hybroid:detection-classifiers}
\end{figure}

Figure~\ref{fig:hybroid:detection-classifiers} separates the learner comparison from the external
baselines. Gradient boosting on the fused representation reaches 0.97 on all three reported
metrics. Random forest reaches 0.96, while the decision tree reaches 0.91. Network flow alone has
a higher point value than program code alone for each learner. Under program code alone, random
forest and gradient boosting both report 0.93 for $F_1$, recall and precision. Gradient boosting
therefore has the highest fused result, as the paper states, while its program-code result ties
with random forest at the published precision \cite{norouzian2021hybroid}.

At two-decimal precision, $F_1$, recall and precision coincide in eight of the nine
configurations. The remaining configuration is gradient boosting over network flow, for which
recall is 0.96 and the other two values are 0.95. Micro-averaging would make the three metrics
equal by construction, and rounding could also hide small differences. The paper does not state
the averaging convention and provides no raw predictions, so neither explanation can be selected
from the surviving evidence \cite{norouzian2021hybroid}.

\section{Malware-Category Classification}
\label{sec:hybroid:10}

The second task assigns each sample to one of five mutually exclusive classes, the four malware
categories of Table~\ref{tab:hybroid:dataset} together with the benign class. The source
publication calls this multi-label classification throughout; it is multi-class classification
with a single label per sample, and this chapter uses the latter term
\cite{norouzian2021hybroid}.

\begin{figure}[t]
  \centering
  \includegraphics[width=\linewidth]{hybroid-categorisation-classifiers.pdf}
  \caption{Five-class malware-category classification by learner and observation modality, drawn
  on the same plan as Figure~\ref{fig:hybroid:detection-classifiers}. The score axis is truncated
  to 0.78--0.96, because this task occupies a lower interval. The values are the five-fold means
  printed to two decimal places in Figure~9 of \cite{norouzian2021hybroid}, checked against the
  bar geometry of that figure. Decision-tree recall is 0.86 for program code and for the fused
  representation, so those two markers are drawn concentrically. No fold-level scores are
  published, so the figure carries no error bars. Redrawn by the candidate.}
  \label{fig:hybroid:categorisation-classifiers}
\end{figure}

Figure~\ref{fig:hybroid:categorisation-classifiers} gives the category-classification results.
Gradient boosting on the fused representation reaches 0.94 on all three metrics. Random forest
follows at 0.92 for $F_1$ and 0.93 for recall and precision. The decision-tree result differs: the
fused representation reaches 0.85 for $F_1$ and precision, below the program-code values of 0.87,
while recall is equal at 0.86. The paper's general statement that combined features give the best
performance therefore has one exception in its own reported figure
\cite{norouzian2021hybroid}.

Table~\ref{tab:hybroid:perclass} gives the area under the ROC curve for each class, for the
gradient boosting classifier on the combined vector. The macro average of 0.976 is reported as a
macro accuracy in the source publication, which its own figure legend contradicts; it is an
average of areas under the curve \cite{norouzian2021hybroid}.

\begin{table}
  \caption{Area under the ROC curve per class, gradient boosting on the combined representation,
  with the pre-exclusion counts of Table~\ref{tab:hybroid:dataset}. Values are read from Figure~10
  of \cite{norouzian2021hybroid}.}
  \label{tab:hybroid:perclass}
  \centering
  \begin{tabular}{lrr}
    \toprule
    Class & Samples & AUC \\
    \midrule
    Benign        & 1,700 & 0.995 \\
    Ransomware    &   101 & 0.991 \\
    Scareware     &   112 & 0.979 \\
    Adware        &   104 & 0.963 \\
    SMS malware   &   109 & 0.946 \\
    \midrule
    Macro average &       & 0.976 \\
    \bottomrule
  \end{tabular}
\end{table}

\begin{figure}[t]
  \centering
  \includegraphics[width=0.72\linewidth]{hybroid-categorisation-roc.pdf}
  \caption{Per-class receiver operating characteristic curves for gradient boosting on the fused
  representation, reproduced from the published vector content of Figure~10 in
  \cite{norouzian2021hybroid}. The legend
  reports a macro-average AUC of 0.976 and the AUC of each class. The paper provides no selected
  operating threshold or fold-level variation.}
  \label{fig:hybroid:categorisation-roc}
\end{figure}

Figure~\ref{fig:hybroid:categorisation-roc} preserves the threshold curves from which the AUC
values in Table~\ref{tab:hybroid:perclass} were read. Benign and ransomware have the largest
reported areas, 0.995 and 0.991. SMS malware has the smallest at 0.946. These values measure class
ranking across thresholds. They do not provide the per-class precision and recall at the decision
rule used for the five-class result, which the paper does not report
\cite{norouzian2021hybroid}.

The averaging convention for the global precision, recall and $F_1$ values is also unstated.
Benign applications remain the majority after the exclusions, so the reported value of 0.94
cannot be interpreted as equal performance across the five classes. Per-class precision and
recall, the fold-level values, and any dispersion measure are unavailable in the paper
\cite{norouzian2021hybroid}.

\section{Modality Comparison}
\label{sec:hybroid:11}

This section answers RQ2. Figures~\ref{fig:hybroid:detection-classifiers}
and~\ref{fig:hybroid:categorisation-classifiers} compare the two single modalities with their
fusion on the same samples, under one protocol, for the three reported learners
\cite{norouzian2021hybroid}.

\begin{table}
  \caption{Difference between the fused score and the better of the two single-modality scores.
  Positive values favour fusion. The differences are computed from Figures~8 and~9 of
  \cite{norouzian2021hybroid}, whose values are printed to two decimal places.}
  \label{tab:hybroid:fusion-effects}
  \centering
  \begin{tabular}{llrrr}
    \toprule
    Task & Learner & $\Delta F_1$ & $\Delta$ recall & $\Delta$ precision \\
    \midrule
    Detection      & Decision tree     & $+0.01$ & $+0.01$ & $+0.01$ \\
    Detection      & Random forest     & $+0.02$ & $+0.02$ & $+0.02$ \\
    Detection      & Gradient boosting & $+0.02$ & $+0.01$ & $+0.02$ \\
    \midrule
    Categorisation & Decision tree     & $-0.02$ & $ 0.00$ & $-0.02$ \\
    Categorisation & Random forest     & $+0.01$ & $+0.01$ & $+0.02$ \\
    Categorisation & Gradient boosting & $+0.03$ & $+0.02$ & $+0.03$ \\
    \bottomrule
  \end{tabular}
\end{table}

The single-modality ordering changes with the task. Network flow has the higher reported $F_1$
for detection under all three learners: 0.90 against 0.89, 0.94 against 0.93, and 0.95 against
0.93. Program code has the higher value for categorisation: 0.87 against 0.80, 0.91 against 0.88,
and 0.91 against 0.90. This reversal is consistent across the three learners reported in the
paper. The experiment does not identify a causal mechanism for it
\cite{norouzian2021hybroid}.

Table~\ref{tab:hybroid:fusion-effects} isolates the point differences that answer RQ2. Fusion adds
one or two points of $F_1$ in binary detection. For categorisation, it reduces decision-tree
$F_1$ by two points and increases the ensemble results by one and three points. Recall for the
decision tree is unchanged. The reported evidence therefore supports a modest fusion improvement
for detection and a learner-dependent result for category classification
\cite{norouzian2021hybroid}.

The paper reports no fold-level values or dispersion measure. The one- and two-point differences
cannot be assigned statistical significance, and the two-decimal rounding prevents a finer
comparison. The chapter therefore reports their magnitude and direction as published point
estimates \cite{norouzian2021hybroid}.

\section{Limitations and Threats to Validity}
\label{sec:hybroid:12}

\paragraph{Encryption, and the age of the corpus.} The strongest limit on external validity is
measurable from the material in Section~\ref{sec:hybroid:7}. Between 0.00 and 10.89 per cent of
the flows in each malware category use TLS, so the evaluation traffic is overwhelmingly
cleartext, and no conclusion about detection over encrypted flows follows from it. The gap has
widened since the capture. Cleartext was permitted by default up to Android 8.1 and has been
disabled by default from Android 9 onwards \cite{aosp2026networksecurity}, and applications
increasingly encrypt content themselves in addition \cite{pourali2022hidden}. A detector fitted
to 2017 traffic would meet a different protocol mix today, and the flow statistics of
Table~\ref{tab:hybroid:flowfeatures} would be observable while the composition of the traffic
they summarise would not be the same.

\paragraph{Part of the flow representation may not describe behaviour.} Two of the four groups in
Table~\ref{tab:hybroid:flowfeatures} are open to this objection. The TCP base sequence numbers
are required to be unpredictable per connection and cannot encode anything about an application.
The TTL and TOS values describe the path and the endpoint rather than the sender. The concern has a
documented precedent. In a widely used intrusion-detection corpus, the time-to-live features
inadvertently fingerprint the operating systems of the hosts involved, and an audit of that
corpus removed them as unjustifiably performant \cite{flood2024bad}. That audit covers different corpora, so the finding does not
transfer, although the mechanism might. Both classes here were captured from the same three
handsets, which leaves little room for the source-side values to separate them, whereas the
destination-side values describe remote endpoints, and command servers and application backends
are different populations of hosts. Whether the mechanism acted here was never measured, and the
feature matrices do not survive, so it stands as an unresolved threat rather than a finding.

\paragraph{The metric convention is unstated.} Section~\ref{sec:hybroid:9} shows that the values
reported for three of the five compared systems cannot be reconciled with a single confusion
matrix on an evaluation cohort whose malicious share lies between 18.2 and 20.5 per cent.
Section~\ref{sec:hybroid:9} also records that precision, recall and $F_1$ coincide at the reported
two-decimal precision in eight of nine detection configurations. An unstated averaging convention
and rounding can both produce this pattern. The paper supplies no raw predictions with which to distinguish them
\cite{norouzian2021hybroid}. The comparison should be read as indicative rather than as a precise
ranking.

\paragraph{The larger experiment is not time-aware.} The 135,905-sample experiment of
Section~\ref{sec:hybroid:4} used a random split. Training and test samples are therefore drawn
from the same period, no temporal ordering is enforced, and no area under the time-decay curve is
reported. Section~\ref{sec:background:evaluation} sets out why an evaluation of this kind
overstates what a detector will achieve against samples that appear after it was trained
\cite{pendlebury2019tesseract}. That experiment is reported here as a check that the pipeline
scales to a larger corpus, and nothing about robustness over time is claimed from it.

\paragraph{Obfuscation.} Graph features could not be extracted from 47 obfuscated applications,
which is 2.2 per cent of the source corpus. The source publication
states in two places that the graph structure is unaffected by obfuscation, which those 47 failures
contradict. The defensible version of the claim is narrower: identifiers and manifest entries can
be renamed without changing the call structure, so a representation built on structure is
insensitive to the transformations that defeat name-based features, and it remains sensitive to
packers and to transformations that rewrite the structure itself.

\paragraph{Two steps are unspecified.} The weights of the method-level mean in
Section~\ref{sec:hybroid:5} are not given, which leaves that mean undefined as published. The step
that composes the per-method vectors into the application-level vector in
Section~\ref{sec:hybroid:6} is not described by any equation. Both are recorded as gaps in the
source rather than reconstructed here.

\paragraph{Fusion is not uniformly beneficial.} Section~\ref{sec:hybroid:11} reports the decision
tree on the category task, where the combined representation scores below the program code alone.
The paper states that combined features give the best performance, while its category figure
reports $F_1$ and precision of 0.85 for the combined decision tree against 0.87 for program code
alone \cite{norouzian2021hybroid}.

\paragraph{The threat model is non-adaptive, and the exposure is structural.} No adversary aware
of this detector is modelled, and no evasion attempt was made against it. That limit has a
specific shape here rather than a generic one. The representation is built over the application's
function-call graph, and attacks published since operate on exactly that object at the level of
the code: semantics-preserving rewriting that injects methods and call edges reaches the input of
the pipeline before any feature is computed, reported at 87.9 per cent attack success against
graph-summary detectors under white-box access \cite{song2026fcghunter}, and a related attack
reports an average of 85.17 per cent across seven detectors whose set includes graph neural
networks, under a perturbation budget of five per cent \cite{li2025efficient}.
Neither was run against \hybroid{}, so nothing here shows it robust and nothing shows it broken.
One asymmetry is worth stating as a hypothesis rather than as a result: the call edges those
operators inject are made unreachable by construction, so they generate no traffic, and the flow
half of the representation would be unchanged by them. Whether that translates into graceful
degradation is untested.

\section{Chapter Summary}
\label{sec:hybroid:13}

This chapter presented \hybroid{}, which detects and categorises Android malware from two
observations of the same application: the structure of its code, recovered without running it and
represented as a function-call graph whose vertices carry embeddings learned from the
instructions of each method and from the control flow between its basic blocks, and the traffic it produces on a physical handset, reduced to \netfeatures{}
features of its bidirectional flows. On CICAndMal2017 the combined representation reaches an
accuracy of 97.0 per cent for detection, and 0.94 on precision, recall and $F_1$ for category
classification, in both cases with a gradient boosting classifier
\cite{norouzian2021hybroid}.

The within-system contrasts in Section~\ref{sec:hybroid:11} answer RQ2. Fusion adds one or two
points of $F_1$ for detection over the better single modality. For category classification the
difference ranges from minus two points for the decision tree to plus three for gradient boosting;
random forest gains one point. Network flow leads program code in detection, while program code
leads network flow in category classification for each reported learner. These are point
differences rounded to two decimals, and the paper reports no fold-level dispersion from which to
test their statistical significance \cite{norouzian2021hybroid}.

Chapter~\ref{chap:hgann} keeps the question of context and changes where the context is sought.
Here it was sought across two sources of evidence, one static and one dynamic, at the cost of
requiring a traffic capture for every sample. There it is sought within the program alone, by
representing relations that involve more than two functions at once, for a detector that needs
only the application package.
